% 原书物理页=印刷页2--3；第1章开头与1.1全节.
\chapter{Lie 代数与 Lie 群}\label{chap:ssl-1}

在本章中，我们先给出\newterm{ssl-lie-algebra}{Lie 代数}{Lie algebra}的基本定义，然后通过说明这些定义如何得自\newterm{ssl-lie-group}{Lie 群}{Lie group}这一自然的观念，来论证它们的合理性.

\section{Lie 代数}\label{sec:ssl-1-1}

设 $\mathbb{K}$ 是一个域.

\begin{definition}\label{def:ssl-1-1-1}
$\mathbb{K}$-Lie 代数是指一个 $\mathbb{K}$-向量空间 $L$，其上配备一个 $\mathbb{K}$-双线性映射
\[
\begin{aligned}
[\;,\;]\colon\ & L\times L\to L,\\
& (X,Y)\mapsto [X,Y],
\end{aligned}
\]
称之为\newterm{ssl-lie-bracket}{Lie 括号}{Lie bracket}，且满足下面两条公理：
\begin{enumerate}[label=\roman*)]
\item 对一切 $X\in L$ 有 $[X,X]=0$（当 $\operatorname{car}(\mathbb{K})\neq 2$ 时，这等价于对一切 $X,Y$ 有 \mbox{$[X,Y]=-[Y,X]$）.}
\item 对一切 $X,Y,Z\in L$ 有 $[X,[Y,Z]]+[Y,[Z,X]]+[Z,[X,Y]]=0$（\newterm{ssl-jacobi-identity}{Jacobi 恒等式}{Jacobi identity}）.
\end{enumerate}
\end{definition}

\begin{remark*}
第二条公理所起的作用，与结合代数定义中的结合律公理相同.
\end{remark*}

\begin{example*}
任何配上零括号的 $\mathbb{K}$-向量空间！
\end{example*}

下面是几个不那么平凡的例子.

\begin{numpara}[结合代数的底层 Lie 代数]\label{np:ssl-1-1-2}
若 $A$ 是一个 $\mathbb{K}$-\newterm{ssl-associative-algebra}{结合代数}{associative algebra}，则 $[x,y]:=xy-yx$ 在 $A$ 上定义了一个 Lie 括号（习题）. 这里最重要的例子是 $\mathbb{K}$-向量空间 $V$ 上全体 $\mathbb{K}$-线性自同态构成的代数 $A=\operatorname{End}_{\mathbb{K}}(V)$. 当 $V$ 有限维时，任取一组基便给出它与 $M_n(\mathbb{K})$ 的一个同构. 由于一个稍后会显现的原因，我们把这样得到的 Lie 代数记作 $\mathfrak{gl}_n(\mathbb{K})$，相应地记作 $\mathfrak{gl}(V)$.
\end{numpara}

\begin{numpara}[代数的导子]\label{np:ssl-1-1-3}
设 $A$ 是一个 $\mathbb{K}$-结合代数. $A$ 的一个\newterm{ssl-derivation}{导子}{derivation}，是指底层 $\mathbb{K}$-向量空间的一个自同态 $\partial$，它满足 Leibniz 公式：$\partial(ab)=(\partial a)b+a(\partial b)$. $A$ 的全体导子之集 $\operatorname{Der}(A)$ 是 $\operatorname{End}_{\mathbb{K}}(A)$ 的一个向量子空间. 两个导子 $\partial,\partial'$ 的复合 $\partial\circ\partial'$ 一般不再是导子（换言之，$\operatorname{Der}(A)$ 不是结合代数 $\operatorname{End}_{\mathbb{K}}(A)$ 的子代数），但 $[\partial,\partial']:=\partial\circ\partial'-\partial'\circ\partial$ 仍是 $A$ 的导子（习题！）. 于是 $\operatorname{Der}(A)$ 是 Lie 代数 $\mathfrak{gl}(A)$ 的一个\newterm{ssl-lie-subalgebra}{Lie 子代数}{Lie subalgebra}.
\end{numpara}

\begin{example*}
设 $M$ 是一个\newterm{ssl-smooth-manifold}{光滑流形}{smooth manifold}（即 $C^\infty$ 类的微分流形）. 我们回顾（可参阅任何一部基础微分几何课程）：所谓\newterm{ssl-vector-field}{向量场}{vector field}，是指 $M$ 的\newterm{ssl-tangent-bundle}{切丛}{tangent bundle} $T_M$ 的光滑截面 $\xi\colon M\longrightarrow T_M$（直观地说，就是对每个 $x\in M$ 给定一个在 $x$ 处切于 $M$ 的向量 $\xi(x)$，且它以 $C^\infty$ 的方式依赖于 $x$）. 这样的向量场定义了 $M$ 上光滑函数代数 $C^\infty(M,\mathbb{R})$ 的一个导子，它把函数 $f$ 送到函数 $x\mapsto\mathrm{d}_x f(\xi(x))$. 通过局部计算可以证明，这样得到的映射 $V(M)\longrightarrow\operatorname{Der}(C^\infty(M,\mathbb{R}))$ 是双射. 由此，$M$ 上全体向量场构成的（无限维）$\mathbb{R}$-向量空间 $V(M)$ 上带有一个 $\mathbb{R}$-Lie 代数结构.
\end{example*}

同样，$\mathbb{K}$-Lie 代数 $L$ 的导子定义为 $\operatorname{End}_{\mathbb{K}}(L)$ 中满足 $\partial([X,Y])=[\partial X,Y]+[X,\partial Y]$ 的元素 $\partial$. 与上面一样，子 $\mathbb{K}$-向量空间 $\operatorname{Der}(L)$ 不关于乘积（复合）封闭，但关于括号封闭，因而是 Lie 代数 $\mathfrak{gl}(L)$ 的 Lie 子代数.

\begin{example*}
下面是 Lie 代数的导子的一个非常重要的例子. 对 $Z\in L$，记
\[
\begin{aligned}
\operatorname{ad}_Z\colon\ & L\to L,\\
& X\mapsto\operatorname{ad}_Z(X):=[Z,X].
\end{aligned}
\]
考虑到括号的反对称性，Jacobi 恒等式恰好告诉我们 $\operatorname{ad}_Z$ 是 $L$ 的导子.
\end{example*}

\begin{definition}\label{def:ssl-1-1-4}
\newterm{ssl-lie-algebra-homomorphism}{Lie 代数同态}{homomorphism of Lie algebras}是指一个与括号交换的 $\mathbb{K}$-线性映射 $f\colon L\longrightarrow L$\footnote{译注：原文此处的目标写作 $L$；按所讨论的一般同态，目标应为 $L^{\prime}$.}，即对一切 $X,Y\in L$，$f([X,Y])=[f(X),f(Y)]$.
\end{definition}

\begin{example*}
迹 $\operatorname{tr}\colon\mathfrak{gl}(V)\longrightarrow\mathbb{K}$ 是一个 Lie 代数同态.（在 $\mathbb{K}$ 上放置唯一可能的 Lie 括号，即零括号.）
\end{example*}

\textbf{特别值得注意的特殊情形：表示.}\ Lie 代数同态 $r\colon L\longrightarrow\mathfrak{gl}_n(\mathbb{K})$ 称为 $L$ 的\newterm{ssl-matrix-representation}{矩阵表示}{matrix representation}. 更内在地说，若 $V$ 是 $\mathbb{K}$-向量空间，则 Lie 代数同态 $r\colon L\longrightarrow\mathfrak{gl}(V)$ 称为 $L$ 的\newterm{ssl-linear-representation}{线性表示}{linear representation}. 这时称 $V$ 为表示 $r$ 的表示空间. 下面是一个非常重大的例子.

\begin{numpara}[Lie 代数的伴随表示]\label{np:ssl-1-1-5}
Jacobi 恒等式表明，映射
\[
\begin{aligned}
\operatorname{ad}\colon\ & L\to\mathfrak{gl}(L),\\
& Z\mapsto\operatorname{ad}_Z
\end{aligned}
\]
是一个 Lie 代数同态，称之为 $L$ 的\newterm{ssl-adjoint-representation}{伴随表示}{adjoint representation}. 如前所见，它的像包含在 $L$ 的导子所构成的 Lie 子代数中；于是 $\operatorname{ad}(L)\subset\operatorname{Der}(L)\subset\operatorname{End}_{\mathbb{K}}(L)$.
\end{numpara}
